\section{Precios de la API: bien común del ecosistema y sustitución comercial}

La API de Reddit atiende a la vez a bots de moderadores, herramientas de investigación, clientes de accesibilidad, proyectos de aficionados, aplicaciones de terceros y usos comerciales de datos a gran escala. Un precio único parece sencillo, pero perjudica a las herramientas de bajo costo y alto valor público; que todo sea gratuito, en cambio, puede dejar a la plataforma con costos de infraestructura cada vez mayores y sin ingresos por publicidad o por licencias de datos. El verdadero problema de diseño es cómo clasificar, medir, notificar y migrar, no solo si «cobrar es correcto».

\subsection{Clasificar a los usuarios por su comportamiento, no por su identidad}

\term{Rate limit} es un límite al número de llamadas por unidad de tiempo, a la concurrencia o al volumen de datos, que sirve para proteger la capacidad y la equidad. El diseño de la API puede asignar un mayor peso de valor público a las herramientas de moderadores y a las aplicaciones de accesibilidad, ofrecer cuotas gratuitas a la investigación y a los proyectos de aficionados, y establecer contratos para los usos comerciales que sustituyen al cliente oficial o extraen datos de forma masiva. Los criterios de clasificación deben ser observables; de lo contrario, el equipo solo puede fijar los precios según la fama de cada empresa y su capacidad de negociación.

\begin{figure}[H]
\centering
\includegraphics[width=0.96\textwidth]{images/08-api-economics.png}
\caption{Niveles de la economía de la API: las herramientas comunitarias, la investigación, los clientes de terceros y el uso comercial a gran escala requieren contratos distintos.}
\figsource{Redibujada a partir de los subtítulos de la entrevista, 32:55--38:00, `images/08-api-economics.png`.}
\end{figure}

\begin{importantbox}{Lectura de la figura: antes del precio, definir la relación de sustitución}
Un moderation bot fortalece a Reddit mismo; un cliente completo de terceros puede sustituir la publicidad y la experiencia del producto oficial; el entrenamiento masivo de modelos traslada el valor del corpus a sistemas externos. Con el mismo número de llamadas, el significado económico puede ser distinto. Un esquema defendible debe hacer públicos las categorías, las cuotas gratuitas, las garantías de rendimiento, las solicitudes de excepción y los periodos de migración, en lugar de tratar a todos los usuarios con un precio unitario único en el último momento.
\end{importantbox}

\subsection{La protesta es retroalimentación, no un veto automático}

Los cambios en la API provocaron el cierre de subreddits y protestas de la comunidad. En la entrevista, Huffman reconoce que las protestas son coherentes con el carácter democrático de Reddit, aunque sostiene que la dirección sigue teniendo que tomar decisiones sobre la sostenibilidad de la plataforma. Ambas cosas pueden ser ciertas a la vez: la comunidad tiene derecho a expresar sus costos colectivos, y la plataforma también debe evaluar la seguridad, los ingresos, la infraestructura y el control del producto a largo plazo. Lo que importa es si el proceso de decisión permite que los afectados entiendan las razones con anticipación, prueben alternativas y dispongan de un plazo de migración razonable.

\begin{knowledgebox}{Nota de clase: no reducir el conflicto a un relato de buenos y malos}
Los desarrolladores de terceros invirtieron años en construir productos, los usuarios desarrollaron dependencia y las herramientas de moderadores asumen trabajo público de la plataforma; Reddit, por su parte, paga los costos de almacenamiento, ancho de banda, seguridad e ingeniería, y enfrenta la sustitución comercial. El conflicto surge de volver a ponerle precio a un antiguo contrato implícito. La revisión posterior más útil no es determinar quién está más enojado, sino identificar qué dependencias no se registraron, qué grupos no tenían una ruta de migración y qué costos nunca fueron transparentes.
\end{knowledgebox}

\begin{table}[H]
\centering
\begin{tabular}{L{0.2\textwidth}L{0.34\textwidth}L{0.36\textwidth}}
\toprule
Etapa & Lo que debe aportar la plataforma & Lo que debe aportar el ecosistema \\
\midrule
Descubrimiento & Clasificación de usos, modelo de costos, inventario de dependencias & Volumen real de llamadas, funciones críticas, necesidades de accesibilidad \\
Diseño & Precios, nivel gratuito, SLA, condiciones de excepción & Alternativas y periodo de migración aceptable \\
Piloto & Entorno aislado de pruebas, monitoreo, condiciones de reversión & Pruebas con poco tráfico y retroalimentación de compatibilidad \\
Migración & Documentación, herramientas, soporte y fechas definidas & Actualización de clientes, avisos a usuarios, exportación de datos \\
Revisión posterior & Interrupciones, quejas, ingresos y pérdidas del ecosistema & Casos fallidos, dependencias heredadas y propuestas de mejora \\
\bottomrule
\end{tabular}
\caption{Un cambio en la política de la API es un proyecto de migración de plataforma, no un anuncio de precios.}
\end{table}

\begin{warningbox}{Las restricciones posteriores a la salida a bolsa cambian el criterio de decisión}
Una empresa que cotiza en bolsa debe rendir cuentas de sus ingresos, costos y riesgos, pero la responsabilidad ante los accionistas no anula automáticamente la responsabilidad ante la comunidad. Los moderadores voluntarios, los creadores de contenido y los desarrolladores de largo plazo aportaron valor a la plataforma, aunque no necesariamente aparezcan en una tabla de costos tradicional. Una gestión madura debe incorporar la salud del ecosistema como indicador observable, en lugar de descubrir las dependencias solo cuando estalla la protesta.
\end{warningbox}

\subsection{Resumen de la sección}

La causa de fondo de la controversia de la API es que distintos patrones de uso compartían una interfaz que históricamente fue casi gratuita. Una solución sostenible requiere clasificación por comportamiento, medición transparente, excepciones por valor público, contratos comerciales y un procedimiento de migración. La apertura no significa precio cero para siempre, y la sostenibilidad tampoco equivale a un cierre repentino; entre ambas hace falta una gobernanza de la interfaz predecible a largo plazo.
